\documentclass[11pt]{article}
\usepackage[letterpaper,margin=0.9in]{geometry}
\usepackage{amsmath,amssymb,amsthm,hyperref}
\hypersetup{hidelinks,pdftitle={Necessary conditions for a large least quadratic nonresidue}}
\newtheorem{proposition}{Proposition}
\newcommand{\LL}{\mathcal L}
\newcommand{\e}{\mathrm e}
\newcommand{\ii}{\mathrm i}
\newcommand{\R}{\mathbb R}
\title{Necessary conditions for a large least quadratic nonresidue}
\author{Research notes}
\date{}
\begin{document}
\maketitle
Assume Dirichlet GRH throughout. Let $p$ run over odd primes,
$\chi=(\cdot/p)$, $n_p=\min\{n\ge1:\chi(n)=-1\}$, and $\LL=\log p$.
The question is whether $n_p=o(\LL^2)$. Failure would give a fixed
$c>0$ and infinitely many primes with $n_p>c\LL^2$. The estimates here
describe what would have to occur at those primes. The argument leaves
their existence open.

\section{A test at the real point $s=1$}
Write $\gamma_E$ for Euler's constant and put
\[
 V=\e^{-\gamma_E}L(1,\chi),\qquad D=-L'(1,\chi)/L(1,\chi).
\]
The Euler product and nonvanishing give $V>0$. All implicit constants
below are absolute unless subscripted.

\begin{proposition}
If $\e\le x<n_p$ and $x<p$, put $H=\log x$ and $u=\LL/\sqrt x$. Then
\begin{align}
 D&=H+O(1+u),\label{d}\\
 \log(V/H)&=-1+(D+\gamma_E)/H+O((1+u)/H^2).\label{joint}
\end{align}
In particular $V\ge H-O(1+u)$.
\end{proposition}

\begin{proof}
Use Lemmas 2.3--2.6 of Lamzouri, Li and Soundararajan \cite{lls}.
Write $B=|\Re B(\chi)|$ for their zero sum, $Q=\tfrac12\log(p/\pi)$,
and $a=(1-\chi(-1))/2$. Lemma 2.3 at $x=4$ gives $B\ll\LL$.
The completed functional equation gives
$D=Q+\tfrac12\psi_0((1+a)/2)-B$, with $\psi_0=\Gamma'/\Gamma$.
At general $x$, the same lemma gives
\[
 (1+O(1/\sqrt x))B=(1-1/x)Q
 -\sum_{n\le x}\frac{\Lambda(n)}{n}(1-n/x)+O(1).
\]
Here $\chi(n)=1$ throughout the sum and their parity terms $E_a(x)$
are bounded for $x\ge\e$. By Lemma 2.4 the sum is
$H-(1+\gamma_E)+O(x^{-1/2})$. These identities give \eqref{d}.
The integrated identity in the proof of Lemma 2.5 gives
\[
 \log L(1,\chi)=
 \sum_{2\le n\le x}\frac{\Lambda(n)}{n\log n}\frac{\log(x/n)}{H}
 +\frac D H+O\!\left(\frac{B}{\sqrt x\,H^2}+\frac{1}{xH^2}\right).
\]
By Lemma 2.6 the sum is $\log H+\gamma_E-1+\gamma_E/H
+O(1/(\sqrt x\,H^2))$. Substitution gives \eqref{joint}. Applying
$\exp z\ge1+z$ to \eqref{joint} and using \eqref{d} proves the lower bound.
\end{proof}

For fixed $c>0$, substituting $x=c\LL^2<n_p$ and expanding the
exponential in \eqref{joint} gives
\[
 V=2\log\LL+O_c(1),\qquad
 V-D=\gamma_E+O_c(1/\log\LL).
\]
There is also a quantitative implication. Suppose for all sufficiently
large primes that $V\le 2\log\LL-h(p)$, where $h(p)\to\infty$.
Take $x=n_p-1$. For $x\ge\e$ the lower bound in Proposition 1 yields
$h(p)\le 2\log u+O(1+u)$, since $H=2\log\LL-2\log u$.
Consequently $u\gg h(p)$ and $n_p\ll\LL^2/h(p)^2$. The bounded
cases also satisfy this estimate because $V>0$ implies $h(p)<2\log\LL$.
Theorem 1.5 of \cite{lls} supplies only a bounded additive saving in $V$.

\paragraph{Test with multiple truncation lengths.}
At the precision of \eqref{d}--\eqref{joint}, all lengths
$\e\le x\le Y=c\LL^2$ can be compatible at once. To see this for
$0<c\le1$, try the scalar values $V_*=\log Y$, $D_*=\log Y-\gamma_E$.
Put $v=\log(Y/x)\ge0$, so $u=\e^{v/2}/\sqrt c$.
The error in \eqref{d} is $v-\gamma_E=O(1+u)$. The absolute error
in \eqref{joint} is at most
\[
 v/H-\log(1+v/H)\le v^2/(2H^2)\le (8/\e^2)\e^{v/2}/H^2 .
\]
Thus repeating these two order estimates cannot force a contradiction.
The scalar values are a consistency test. Realizing them by an actual
quadratic character would require arithmetic information.

\section{A count of residues at a longer length}
Set $z=\LL^2$ and $X=z^2$. Granville and Soundararajan give
$1-\log 2=\rho(2)$ as the GRH upper limit for the normalized short
character sum at this length, even allowing all fundamental discriminants
\cite[Theorem 2 and the discussion of $\alpha(B)$]{gs}. For the proposed
counterexamples one can control the distance from this limit.

\begin{proposition}
For fixed $c>0$, at every sufficiently large prime with $n_p>cz$,
\[
 \frac{1}{X}\sum_{n\le X}\chi(n)
 =1-\log 2+O_c(1/\log z).
\]
In particular a uniform upper bound $1-\log 2-h(p)/\log\log p$ on
the left, for any $h(p)\to\infty$, implies $n_p=o(\LL^2)$.
\end{proposition}

\begin{proof}
First, for every $t\ge\e$, without any restriction on $n_p$, the same
two formulas used for Proposition 1 give
\[
 \log L(1,\chi)
 =\sum_{2\le n\le t}\frac{\Lambda(n)\chi(n)}{n\log n}
 +O\!\left(\frac{1+\LL/\sqrt t}{\log t}\right).
\]
Indeed Lemma 2.3 retains $\chi(n)$ in the prime sum, giving
$D-\sum_{n\le t}\Lambda(n)\chi(n)/n=O(1+\LL/\sqrt t)$.
Removing the weight $1-n/t$ here costs at most
$t^{-1}\sum_{n\le t}\Lambda(n)=O(1)$. In the integrated formula from
Lemma 2.5 this difference is precisely the correction from removing
the logarithmic weight, after division by $\log t$. Subtract the
approximations at $t=z$ and $t=X$. Higher prime powers exceeding $z$
contribute $O(z^{-1/2})$ in absolute value. This follows by summing the
geometric tails separately for primes below and above $\sqrt z$. Hence
\[
 \sum_{z<\ell\le X,\ \ell\ {\rm prime}}\frac{\chi(\ell)}{\ell}
 =O(1/\log z).
\]

Temporarily replace $\chi(\ell)$ by $+1$ for primes $\ell\le z$ and
extend completely multiplicatively to obtain $f$. Since $\chi(\ell)=1$
below $cz$ and $X<p$, this changes the sum over $n\le X$ by at most
$2X\sum_{cz<\ell\le z}1/\ell=O_c(X/\log z)$, with the sum over primes.
Each integer up to $X=z^2$ has at most one prime factor exceeding $z$,
counting multiplicity. It follows exactly that
\[
 \sum_{n\le X} f(n)=\lfloor X\rfloor-
 \sum_{z<\ell\le X,\ \ell\ {\rm prime}}(1-\chi(\ell))\lfloor X/\ell\rfloor .
\]
The total floor error is $O(\pi(X))=O(X/\log z)$. By Mertens' estimate
the prime reciprocal sum from $z$ to $z^2$ is $\log 2+O(1/\log z)$.
This proves the formula. Failure of the little-$o$ claim supplies a
fixed $c$ and an infinite subsequence, contradicting the proposed upper bound.
\end{proof}

In ordinary counts the proportion of residues would therefore be
$1-\tfrac12\log 2+O_c(1/\log\log p)$, about $65.34\%$.
For comparison, the error in the smooth-number cutoff in Theorem 2
of \cite{gs} permits the limiting endpoint. Its conclusion alone gives
no uniform downward saving at the scale in Proposition 2.

\subsection*{The same prime at several lengths}
Write $P^+(n)$ for the largest prime factor of $n$, with $P^+(1)=1$,
and $\Psi(X,z)=\#\{n\le X:P^+(n)\le z\}$. One can keep the same
error uniformly over all lengths. The error below is additive and
can exceed the main term when $X$ is large.

\begin{proposition}
With $z=(\log p)^2$, uniformly for sufficiently large $p$ and $X\ge z$,
\[
 \sum_{n\le X}\chi(n)
 =\sum_{\substack{n\le X\\ P^+(n)\le z}}\chi(n)+O(X/\log z).
\]
If $n_p>cz$ for fixed $c>0$, the right side equals
$\Psi(X,z)+O_c(X/\log z)$.
\end{proposition}

\begin{proof}
Put $S(t)=\sum_{n\le t}\Lambda(n)\chi(n)\log(t/n)$.
Lemma 2.2 of \cite{lls} and $|\Re B(\chi)|\ll\log p$ give
$S(t)\ll\sqrt t\log p$ for $t\ge 2$. For $t\ge z$ take
$\eta=(z/t)^{1/4}\le1$. The difference
$(S(t\e^\eta)-S(t))/\eta$ equals
$\sum_{n\le t}\Lambda(n)\chi(n)$ up to an error at most
$\psi(t\e^\eta)-\psi(t)$, where $\psi$ is the untwisted Chebyshev
function. The RH estimate $\psi(t)=t+O(\sqrt t\log^2(2t))$ gives
an error $O(t\eta)$ since $\eta\ge t^{-1/4}$. The difference itself
is $O(\sqrt t\log p/\eta)=O(t\eta)$. Thus
\[
 \left|\sum_{n\le t}\Lambda(n)\chi(n)\right|\ll z^{1/4}t^{3/4}.
\]
For $2\le t<z$ the trivial Chebyshev bound implies the same inequality.
Removing the proper prime powers costs $O(\sqrt t\log^2(2t))$.
Partial summation now gives, for $t\ge z$,
\[
 \left|\sum_{\ell\le t,\ \ell\ {\rm prime}}\chi(\ell)\right|
 \ll z^{1/4}t^{3/4}/\log t .
\]

For every $n\le X$ with $P^+(n)>z$, write $n=m\ell$ with
$\ell=P^+(n)$. For fixed $m\le X/z$ the admissible primes satisfy
$z<\ell\le X/m$ and $\ell\ge P^+(m)$. The last bound at the two
ends of this interval gives $O(z^{1/4}(X/m)^{3/4}/\log z)$ for
their signed sum, including a possible lower endpoint. Summing over
$m$ uses $\sum_{m\le X/z}m^{-3/4}\ll(X/z)^{1/4}$ and proves the
first assertion. Under $n_p>cz$, every $z$-smooth integer with negative
sign has a prime factor in $(cz,z]$. A union bound over multiples and
the prime reciprocal estimate from Proposition 2 prove the second.
\end{proof}

In particular, if there were an exceptional sequence of primes, then
for every fixed $u\ge1$ the residue fraction through $z^u$ on that
same sequence would tend to $(1+\rho(u))/2$. Here $\rho$ is the
Dickman function, the limiting density of $X^{1/u}$-smooth integers
\cite[Introduction]{gs}. For $1\le u\le2$ the explicit fraction is
$1-\tfrac12\log u+O_c(1/\log z)$, uniformly in $u$, by counting
multiples of primes in $(z,z^u]$ and applying Mertens' estimate.
At $u=3$ the limiting fraction is about $52.43\%$. An upper bound
strictly below one such limit for every sufficiently large prime
would rule out the sequence.

\paragraph{Can multiplicativity rule out these counts jointly?}
Fix any finite set of lengths $X_j\ge z$. Set $f(\ell)=1$ for primes
$\ell\le z$ and choose independent fair signs for the remaining primes,
extending $f$ completely multiplicatively. Write $S_f(X)=\sum_{n\le X}f(n)$.
The expectation of $f(n)$ vanishes unless every prime exceeding $z$
occurs in $n$ to even order. Beyond the smooth integers this requires
$\ell^2\mid n$ for some prime $\ell>z$. Consequently
\[
 0\le \mathbb E S_f(X)-\Psi(X,z)
 \le X\sum_{\ell>z}\ell^{-2}\ll X/z ,
\]
where the sum is over primes. Changing one sign changes $S_f(X)$ by
at most $2X/\ell$. Exposing the independent signs one at a time and
adding the variances of the martingale differences therefore gives
$\operatorname{Var} S_f(X)\ll X^2/z$. Chebyshev's inequality and a
union bound give a single choice of signs with
$S_f(X_j)=\Psi(X_j,z)+O_k(X_j/\sqrt z)$ at all $k$ prescribed lengths.
Thus complete multiplicativity alone can reproduce all these counts
to smaller error than Proposition 3 requires.

In fact any fixed such sign pattern through $\max_j X_j$ occurs
for infinitely many prime moduli. Choose $p\equiv1\pmod8$ and, for
each odd prime $\ell$ in range, prescribe a nonzero class for $p$
modulo $\ell$ with the desired Legendre symbol. This handles $2$
since $z\ge2$ and $f(2)=1$. The Chinese remainder theorem, Dirichlet's
theorem, and quadratic reciprocity then give the claim. The missing
constraint is the size of the modulus. A prime obtained this way may
have $(\log p)^2$ much larger than the cutoff $z$ used to design the signs.

\section{A test on a bounded interval of zero heights}
For the zeros $\tfrac12+\ii t_j$ of $L(s,\chi)$, with multiplicity,
define a nonnegative mean square
\[
 F_p(y,T)=\int_{\R}\left|\sum_{|t_j|\le T}\exp(\ii t_j(\log y+v))\right|^2
                  \e^{-2|v|}\,dv .
\]
This uses both signs of the ordinates. Kandhil, Languasco and Moree
use a corresponding positive-ordinate sum \cite[Section 4]{klm}.

\begin{proposition}
For each $c>0$ there are fixed $T_c$ and $b_c>0$ such that, at all
sufficiently large primes with $n_p>y=c\LL^2$,
\[
 F_p(y,T_c)\ge b_c\LL^2 .
\]
In particular, $F_p(c\LL^2,T)=o_{c,T}(\LL^2)$ for each fixed $c,T>0$
would imply the desired growth rate.
\end{proposition}

\begin{proof}
Fix a real, nonnegative, nonzero $g\in C_c^\infty((-1,0))$ and put
$\Phi(t)=\int g(v)\e^{\ii tv}\,dv$ and $A=\int g(v)\e^{v/2}\,dv>0$.
The smooth explicit formula for a primitive nonprincipal character gives
\[
 \sum_{n\ge2}\frac{\Lambda(n)\chi(n)}{\sqrt n}g(\log(n/y))
 =-\sum_j\e^{\ii t_j\log y}\Phi(t_j)+O_g(y^{-1/2}).
\]
Indeed, apply Mellin inversion to $-L'/L$ with transform
$y^{s-1/2}\int g(v)\e^{(s-1/2)v}\,dv$ and shift left past its zeros.
The trivial zeros at $-a-2k$ contribute $O_g(y^{-1/2})$ by compact
support. The prime sum equals $\sqrt y(A+o_g(1))$ by the prime number
theorem, since all its arguments are below $y$.

The standard unit-interval zero count is $O(\log(p(|t|+2)))$
uniformly in $p,t$ by \cite[Section 3]{fourier}. Rapid decay of $\Phi$
therefore makes the absolute sum over $|t_j|>T$ at most
$A\sqrt c\,\LL/4$ for suitable fixed $T=T_c$ and large $p$.
The remaining sum has magnitude at least $A\sqrt c\,\LL/2$.
It is the integral against $g(v)$ of the finite exponential sum defining
$F_p(y,T)$. Weighted Cauchy--Schwarz bounds its square by
$F_p(y,T)\int |g(v)|^2\e^{2|v|}\,dv$, as required.
\end{proof}

\paragraph{Test of zero counts at this frequency.}
The size of the error in the count of ordinates matters. For example,
fix $0<c<1/16$, set $\lambda=\log(c\LL^2)$ and consider the even,
positive model density on $\R$
\[
 d\nu(t)=\frac{\LL}{2\pi}\left(1+2\Re(a(t)\e^{-\ii\lambda t})\right)dt,
 \qquad a(t)=\frac{-\sqrt c}{1/2-\ii t}.
\]
Positivity follows from $2|a|\le4\sqrt c<1$. Integration by parts,
using $\|a'\|_1<\infty$, bounds its discrepancy from
$\LL\,dt/(2\pi)$ on every interval by $O_c(\LL/\lambda)$. This has the
same order as Selberg's fixed-height bound for the actual count discrepancy.
Yet for every $g\in C_c^\infty((-1,0))$ as above the
Riemann--Lebesgue lemma and Fourier inversion give
\[
 \frac1{\LL}\int \e^{\ii\lambda t}\Phi(t)\,d\nu(t)
 \longrightarrow \frac1{2\pi}\int \Phi(t)a(t)\,dt
 =-\sqrt c\int g(v)\e^{v/2}\,dv .
\]
The last equality uses the Fourier transform of
$\mathbf 1_{v<0}\e^{v/2}$. One can replace the density on the positive
axis by a point at each half-integer quantile of its cumulative mass and
reflect these points about zero. This changes interval counts by $O(1)$.
By integration by parts it changes the last normalized integral by
$O_g((1+\lambda)/\LL)=o(1)$. Thus even simple, symmetric model ordinates
can meet both the fixed-height count bound and the smoothed prime-side
constraint at this scale. The model imposes neither an Euler product
nor the full explicit formula at other scales. Ruling out the phase
alignment for actual characters remains the substantive missing estimate.

\begin{thebibliography}{9}\small
\bibitem{lls} Y.~Lamzouri, X.~Li and K.~Soundararajan,
\emph{Conditional bounds for the least quadratic non-residue and related problems},
\href{https://arxiv.org/abs/1309.3595}{arXiv:1309.3595}, version 3.
\bibitem{gs} A.~Granville and K.~Soundararajan, \emph{Large character sums},
\href{https://arxiv.org/abs/math/9903196}{arXiv:math/9903196}.
\bibitem{fourier} E.~Carneiro, M.~Milinovich, E.~Quesada-Herrera and
A.~Ramos, \emph{Fourier optimization, the least quadratic non-residue,
and the least prime in an arithmetic progression},
\href{https://arxiv.org/abs/2404.08380}{arXiv:2404.08380}.
\bibitem{klm} N.~Kandhil, A.~Languasco and P.~Moree,
\emph{Beyond the Riemann Hypothesis bounds: A pair-correlation approach to
the least prime in arithmetic progression and the smallest quadratic non-residue},
\href{https://arxiv.org/abs/2607.14515}{arXiv:2607.14515}.
\end{thebibliography}
\end{document}
